% Mechanization map for the appendix:
% def:representation: AbstractMerge.Representation; GuardedRawJoin.Unique.
% def:metadata: MetadataDependencies; MetadataDependencies.Closed/Past;
%   MetadataDependencies.closed_diff_of_max.
% def:merge-vcs: AbstractMRDT.Raw.Context and MergeVCs in GuardedRawJoin.lean.
% def:replay-supply: AbstractMRDT.Raw.PeelChoice and ReplaySupply.
% thm:join: AbstractMRDT.Raw.join_at_sizes and representationJoin_of_vcs.
% thm:execution: CertifiedClosedExecution.ClosedJoinAt,
%   canonicalConfig_of_execution and virtualMergeBaseState_canonical;
%   CertifiedRGAExecution.Embedded.represented_of_canonical,
%   canonical_of_represented and closedJoinAt are a concrete adapter example.

\section{Concrete representations and merge induction}

The merge argument must retain enough information to reconstruct an event's
metadata, including commuting predecessors when necessary. The following
contract states the concrete equations used by the checked merge induction.
It establishes representation preservation. An independent sequential-history
argument is still required for the correctness criterion.

Write $u_e(s)$ for the update of state $s$ by the original event $e$, and
$m(l,a,b)$ for three-way merge. In this section only, $a\leftrightarrow b$
means commutation on \emph{all} concrete states. For an operation policy
$P$ on update payloads, define the raw semantic order
\[
\begin{aligned}
 \ell^C_H(a,b)\iff{}&
 (a\mathrel{\prec_C}b\land\neg(a\leftrightarrow b))\\
 &\lor\bigl(\neg(a\mathrel{\prec_C}b)
 \land\neg(b\mathrel{\prec_C}a)\land P(a.o,b.o)\\
 &\hspace{12mm}\land\neg\exists c\in H.
 b\mathrel{\prec_C}c\land\neg(b\leftrightarrow c)\bigr).
\end{aligned}
\]
These are the relations consumed by the raw merge theorem; they do not
replace invariant-scoped commutation in the final history criterion.

\begin{definition}[Indexed concrete representation]\label{def:representation}
A representation is a predicate
\[
 R:\mathsf{ReplayContext}\times\mathcal P(\mathcal E)\times\Sigma
       \longrightarrow\mathsf{Prop}.
\]
The assertion $R(C,H,s)$ relates the \emph{entire concrete state} $s$ to a
specific event set $H$ in context $C$. It need not be defined as reachability,
and it is not inferred from the desired query answer. Define
\[
\begin{aligned}
 \mathsf{Supp}_C(H)&\iff H\subseteq\mathcal E_C,\\
 \mathsf{Unique}(R)&\iff
   \forall C,H,s,t.\ R(C,H,s)\land R(C,H,t)\Longrightarrow s=t.
\end{aligned}
\]
An enumeration $\pi\equiv H$ is a finite list without repetitions whose
members are exactly $H$. Thus every finite-set assumption below is an
explicit enumeration assumption.
\end{definition}

\begin{definition}[Metadata dependencies]\label{def:metadata}
For every context $C$, choose a relation $M_C$ on events satisfying
\[
\begin{aligned}
 M_C(a,b)&\Longrightarrow a\mathrel{\prec_C}b,\\
 a\mathrel{\prec_C}b\land\neg(a\leftrightarrow b)
    &\Longrightarrow M_C(a,b).
\end{aligned}
\]
It may retain additional visible predecessors which commute. Define
\[
\begin{aligned}
 \mathsf{Closed}_{M_C}(H)&\iff
   \forall a,b.\ M_C(a,b)\land b\in H\Longrightarrow a\in H,\\
 \mathsf{Past}_{M_C}(e)&=\{x\mid x=e\ \lor\ M_C^+(x,e)\},\\
 \mathsf{Max}^{\ell}_{C,H}(e)&\iff
   \forall x\in H\setminus\{e\}.\ \neg\ell^C_H(e,x),\\
 \mathsf{Max}^{M}_{C,H}(e)&\iff
   \forall x\in H\setminus\{e\}.\ \neg M_C(e,x).
\end{aligned}
\]
The two maximality conditions are separate. If $H\subseteq U$ is
$M_C$-closed and $e$ is $M_C$-maximal in $U$, then
$H\setminus\{e\}$ remains $M_C$-closed. Semantic maximality alone does not
justify this step. For the embedded sequence and anchored queue ports,
$M_C=\prec_C$; their reconstruction past retains all causal metadata.
\end{definition}

\begin{definition}[Five concrete merge VCs]\label{def:merge-vcs}
The following conditions are universally quantified over the displayed
contexts, histories, events and states. Define the common hypotheses
\[
\begin{aligned}
 A_C(H)&:=\mathsf{Supp}_C(H)\land\mathsf{Closed}_{M_C}(H),\\
 K_C(H_1,H_2,e)&:=
   \mathsf{Trans}(\prec_C)\land\mathsf{Irrefl}(\prec_C)
   \land A_C(H_1)\land A_C(H_2)\\
 &\hspace{12mm}\land\mathsf{Max}^{\ell}_{C,H_1\cup H_2}(e)
   \land\mathsf{Max}^{M}_{C,H_1\cup H_2}(e).
\end{aligned}
\]
Within each condition, abbreviate $R_C(H,s):=R(C,H,s)$,
$D:=\mathsf{Past}_{M_C}(e)$ and $H^-:=H\setminus\{e\}$.
These abbreviations compress only repeated hypotheses.
\begin{enumerate}[label=\textbf{VC\arabic*.},leftmargin=*,itemsep=.65em]
\item \emph{Merge commutativity.} From
\[
 A_C(H_1),\ A_C(H_2),\ R_C(H_1\cap H_2,l),\ R_C(H_1,a),\ R_C(H_2,b)
\]
conclude $m(l,a,b)=m(l,b,a)$. This equation is required on represented
triples, not on arbitrary concrete states.

\item \emph{Initial branch.} From $A_C(H)$ and $R_C(H,s)$ conclude
\[
 m(\sigma_0,\sigma_0,s)=s.
\]

\item \emph{Local redistribution.} Assume $K_C(H_1,H_2,e)$,
$e\in H_1\setminus H_2$, and \emph{all} the representation premises
\[
\begin{gathered}
 R_C(H_1\cap H_2,l),\quad R_C(D^-,B),\quad
 R_C(H_1^-,t),\quad R_C(H_2,b),\\
 R_C(D,u_e(B)),\quad R_C(H_1,m(B,t,u_e(B))),\\
 R_C((H_1\cup H_2)^-,m(l,t,b)).
\end{gathered}
\]
Then
\[
 m(l,m(B,t,u_e(B)),b)=m(B,m(l,t,b),u_e(B)).
\]

\item \emph{Shared redistribution.} Assume $K_C(H_1,H_2,e)$,
$e\in H_1\cap H_2$, and
\[
\begin{gathered}
 R_C((H_1\cap H_2)^-,t_0),\quad R_C(D^-,B),\\
 R_C(H_1^-,t_1),\quad R_C(H_2^-,t_2),\quad R_C(D,u_e(B)),\\
 R_C(H_1\cap H_2,m(B,t_0,u_e(B))),\\
 R_C(H_1,m(B,t_1,u_e(B))),\quad
 R_C(H_2,m(B,t_2,u_e(B))),\\
 R_C((H_1\cup H_2)^-,m(t_0,t_1,t_2)).
\end{gathered}
\]
Then
\[
\begin{aligned}
 &m\bigl(m(B,t_0,u_e(B)),m(B,t_1,u_e(B)),m(B,t_2,u_e(B))\bigr)\\
 &\hspace{20mm}=m\bigl(B,m(t_0,t_1,t_2),u_e(B)\bigr).
\end{aligned}
\]

\item \emph{Causal delta.} Assume transitive, irreflexive visibility,
$A_C(U)$, $e\in U$, both $\mathsf{Max}^{\ell}_{C,U}(e)$ and
$\mathsf{Max}^{M}_{C,U}(e)$, and
\[
 R_C(U^-,s),\quad R_C(D^-,B),\quad R_C(D,u_e(B)),\quad R_C(U,u_e(s)).
\]
Then $m(B,s,u_e(B))=u_e(s)$.
\end{enumerate}
The represented smaller-union merge in VC3 and VC4 is obtained by a
\emph{strictly smaller} induction call. Its presence is not a premise
asserting the desired join theorem for the original union. The represented
side reconstructions are also derived on proper subsets, or by VC5 when a
side equals the whole union.
\end{definition}

\begin{definition}[Replay supply and a peel witness]\label{def:replay-supply}
At a fixed context $C$, replay supply consists of both clauses:
\begin{enumerate}[label=(\roman*),leftmargin=*,itemsep=.25em]
\item For every $H$ and $\pi$, if $\pi\equiv H$ and
$\mathsf{Supp}_C(H)$, some $s$ satisfies $R_C(H,s)$.
\item For every represented, supported, nonempty, $M_C$-closed $U$, there
exist $e\in U$ and states $s,B$ satisfying
\[
\begin{gathered}
 \mathsf{Max}^{\ell}_{C,U}(e),\qquad\mathsf{Max}^{M}_{C,U}(e),\\
 R_C(U^-,s),\quad R_C(\mathsf{Past}_{M_C}(e)^-,B),\\
 R_C(\mathsf{Past}_{M_C}(e),u_e(B)),\quad R_C(U,u_e(s)).
\end{gathered}
\]
\end{enumerate}
The second clause is the \emph{peel witness}. It supplies update
reconstruction, not a merge-preservation assumption. Its existence is a
separate datatype obligation; acyclicity of $\ell$ alone supplies neither
metadata maximality nor the four representation assertions.
\end{definition}

\begin{theorem}[Finite concrete join induction]\label{thm:join}
Assume the five VCs, $\mathsf{Unique}(R)$, and
\[
\begin{aligned}
 &\forall C,H,s.\ R(C,H,s)\Longrightarrow R(C,\varnothing,\sigma_0),\\
 &\forall C,H,s.\ R(C,H,s)\Longrightarrow\exists\pi.\ \pi\equiv H.
\end{aligned}
\]
Fix a context $C$ with transitive, irreflexive visibility and replay supply.
For all $H_1,H_2,l,a,b$, if $A_C(H_1)$, $A_C(H_2)$ and
\[
 R_C(H_1\cap H_2,l),\qquad R_C(H_1,a),\qquad R_C(H_2,b),
\]
then
\[
 R_C(H_1\cup H_2,m(l,a,b)).
\]
\end{theorem}
\begin{proof}[Proof structure]
Induct strongly on the length of a duplicate-free enumeration of
$H_1\cup H_2$. An empty side uses uniqueness, the initial-representation
hypothesis and VC2, with VC1 if the right side is empty. Otherwise replay
supply chooses a peel witness. Removing its event preserves metadata closure
by $M_C$-maximality. Supply gives represented predecessor states for the
branches and their intersection. The smaller induction calls establish
side reconstruction and the smaller-union merge representation. Uniqueness
identifies the reconstructed sides with the original sides; VC3 or VC4
then brings the event outside. VC5 and the peel witness's represented update
reconstruct the union. Every recursive call has a strictly smaller event
set; none assumes representation preservation of an unresolved merge.
\end{proof}

\begin{remark}[The full-causal-closure join contract]
The checked derived contract quantifies over fully causally closed
$H_1,H_2$, not over all merely conflict-closed histories. To obtain it from
Theorem~\ref{thm:join}, require replay supply for every $C,H_1,H_2,a,b$
with transitive, irreflexive visibility, supported sides and
$R_C(H_1,a),R_C(H_2,b)$. This is the quantifier order of the supply hypothesis
in the checked derived theorem. Causal closure implies $M_C$-closure because
$M_C\subseteq\prec_C$. Hence the theorem gives the concrete
\emph{representation join} contract
\[
\begin{gathered}
 \mathsf{Supp}_C(H_1),\ \mathsf{Supp}_C(H_2),\quad
 \mathsf{CausalClosed}_C(H_1),\ \mathsf{CausalClosed}_C(H_2),\\
 R_C(H_1\cap H_2,l),\ R_C(H_1,a),\ R_C(H_2,b)
 \quad\Longrightarrow\quad R_C(H_1\cup H_2,m(l,a,b)),
\end{gathered}
\]
under transitive, irreflexive visibility. This contract includes intermediate
representation evidence which canonical query agreement by itself does not
provide.
\end{remark}

\begin{theorem}[Lifting to certified executions]\label{thm:execution}
Fix an internal replay policy whose directed event relation is $\rho$.
This auxiliary policy supplies concrete replay witnesses; it is separate from
the public payload policy and from the revised invariant-scoped history order.
Let $K_\rho(a,b):=\rho(a,b)\lor\rho(b,a)$. Its set-relative order is
\[
\begin{aligned}
 \lambda^\rho_{C,H}(a,b)\iff{}&
 (a\prec_C b\land K_\rho(a,b))\\
 &\lor\bigl(\neg(a\prec_C b)\land\neg(b\prec_C a)\land\rho(a,b)\\
 &\hspace{12mm}\land\neg\exists c\in H.\ b\prec_C c\land K_\rho(b,c)\bigr).
\end{aligned}
\]
Define
\[
 \mathsf{Can}^{\rho}_C(H,s)\iff
 \exists\pi.\ \pi\equiv H\land
 \mathsf{Respects}(\pi,\lambda^\rho_{C,H})\land
 u^*(\sigma_0,\pi)=s.
\]
Suppose, for \emph{every} configuration $C$ satisfying
original mint honesty, the following full-causal-closure canonical join
contract holds:
\[
\begin{gathered}
 \mathsf{Trans}(\prec_C),\ \mathsf{Irrefl}(\prec_C),\quad
 \mathsf{Supp}_C(H_1),\ \mathsf{Supp}_C(H_2),\\
 \mathsf{CausalClosed}_C(H_1),\ \mathsf{CausalClosed}_C(H_2),\\
 \mathsf{Can}^{\rho}_C(H_1\cap H_2,l),\quad
 \mathsf{Can}^{\rho}_C(H_1,a),\quad\mathsf{Can}^{\rho}_C(H_2,b)\\
 \Longrightarrow\quad\mathsf{Can}^{\rho}_C(H_1\cup H_2,m(l,a,b)).
\end{gathered}
\]
Every ordinary or recursive-virtual issuance-certified execution from the
initial configuration then has $\mathsf{Can}^{\rho}_C(H,s)$ at every allocated
version $C.S(v)=(s,H)$. Its event sets are supported and causally closed;
visibility is transitive and irreflexive. Moreover, if the store invariant and
this canonical configuration invariant hold, its recursive virtual base for
versions $(s_1,H_1)$ and $(s_2,H_2)$ is canonical for $H_1\cap H_2$.
\end{theorem}

The displayed join hypothesis is derived rather than supplied as a datatype
axiom: the five VCs, uniqueness, initial representation, finite enumeration
and replay supply yield representation join. A port additionally proves
\[
 R_C(H,s)\quad\Longleftrightarrow\quad\mathsf{Can}^{\rho}_C(H,s)
\]
for supported, causally closed histories in each mint-honest context
with transitive, irreflexive visibility. The backward implication makes the
three canonical inputs represented; representation join constructs the union;
the forward implication makes its result canonical. The embedded ports have
exactly this correspondence. Other ports can use stronger represented
canonical-state certificates rather than asserting this equivalence generally.

The execution induction carries the store invariant together with canonicality.
Ordinary merges use the actual greatest common ancestor's intersection history.
This event-intersection coherence is an admissibility property of the
configuration in Definition~\ref{def:execution}; for operational stores it is
proved from the store invariant, rather than from graph ancestry alone.
Recursive virtual merges use a nested induction on finite ancestry support:
the inner fold represents the union of its inputs, and each recursively
computed base represents their intersection. Original issuance guards remain
at original mint states. Neither intermediate peeling nor scratch computation
requires reissuing an event. This theorem is a merge-and-replay guarantee;
legality and specification visibility of one independent sequential word are
additional obligations of the final correctness bridge.
